\clearpage
\section{Key Specifications}
\label{sec:summary}
% generated by FR3_Characterization/plot_datasheet_spurs.py -- do not edit
\newcommand{\spurCount}{50}
\newcommand{\spurLoMin}{7}
\newcommand{\spurLoMax}{13.125}
\newcommand{\spurLoStepMHz}{125}
\newcommand{\spurTwoLoMin}{14}
\newcommand{\spurTwoLoMax}{26.25}
\newcommand{\spurTwoLoLimit}{26.5}
\newcommand{\spurLoLimit}{13.25}
\newcommand{\spurMax}{−22.5}
\newcommand{\spurMaxLo}{8.125}
\newcommand{\spurMaxTwoLo}{16.25}
\newcommand{\spurMedian}{−32.9}
\newcommand{\spurMin}{−49.8}
\newcommand{\spurMinLo}{10.875}
\newcommand{\spurMinTwoLo}{21.75}
\newcommand{\spurExpected}{−20}
\newcommand{\spurLmxMinC}{58}
\newcommand{\spurLmxMaxC}{71}
\newcommand{\spurToneCount}{25}
\newcommand{\spurToneWithinOne}{21}
\newcommand{\spurToneMedianDiff}{0.4}
\newcommand{\spurAmpLow}{−14.7}
\newcommand{\spurAmpHigh}{−10.8}
\newcommand{\spurZeroBiasWorst}{−13.2}
\newcommand{\spurZeroBiasLoMax}{10}
\newcommand{\rxLeakLoCount}{156}
\newcommand{\rxLeakLoMin}{7}
\newcommand{\rxLeakLoMax}{−49.3}
\newcommand{\rxLeakLoMedian}{−59.6}
\newcommand{\rxLeakLoMaxAt}{17.1}
\newcommand{\rxLeakLoBest}{−85.2}
\newcommand{\rxLeakTwoLoCount}{63}
\newcommand{\rxLeakTwoLoMax}{−58.8}
\newcommand{\rxLeakTwoLoMedian}{−71.1}
\newcommand{\rxLeakTwoLoMaxAt}{13.2}
\newcommand{\rxLeakLoFineMax}{13.2}
\newcommand{\rxLeakLmxMinC}{77}
\newcommand{\rxLeakLmxMaxC}{84}
\newcommand{\rxLeakLoTop}{22.5}
\newcommand{\rxLeakLoBestAt}{7}
\newcommand{\rxLeakOffsetMv}{5}
\newcommand{\rxLeakOffsetCount}{4}
\newcommand{\rxLeakOffsetMaxChange}{0.3}
\newcommand{\rxLeakOffsetLos}{8, 12, 16, 20}
% the 2xLO spur row and note quote these (the spur section itself comes later)
% generated by FR3_Characterization/make_rx_conditions.py -- do not edit
\providecommand{\attn}[1]{\textcolor{red}{\textbf{#1}}}
\newcommand{\condRxGainMinC}{69}
\newcommand{\condRxGainMaxC}{75}
\newcommand{\condRxFlatMinC}{70}
\newcommand{\condRxFlatMaxC}{77}
\newcommand{\condRxCompMinC}{72}
\newcommand{\condRxCompMaxC}{81}
\newcommand{\condRxImageMinC}{75}
\newcommand{\condRxImageMaxC}{82}
\newcommand{\condRxNfMinC}{77}
\newcommand{\condRxNfMaxC}{84}
%
% generated by FR3_Characterization/plot_datasheet_image.py from image_rx_20261009 -- do not edit
\newcommand{\imgLowBypassMedian}{−1}
\newcommand{\imgLowBypassPten}{−4}
\newcommand{\imgLowBypassMin}{−6}
\newcommand{\imgLowBypassCount}{209}
\newcommand{\imgOutlierAX}{0.6962}
\newcommand{\imgOutlierAY}{0.3328}
\newcommand{\imgOutlierBX}{0.8057}
\newcommand{\imgOutlierBY}{0.3857}
\newcommand{\imgLowOptimalMedian}{51}
\newcommand{\imgLowOptimalPten}{24}
\newcommand{\imgLowOptimalMin}{0}
\newcommand{\imgLowOptimalCount}{209}
\newcommand{\imgHighBypassMedian}{1}
\newcommand{\imgHighBypassPten}{−2}
\newcommand{\imgHighBypassMin}{−6}
\newcommand{\imgHighBypassCount}{203}
\newcommand{\imgHighOptimalMedian}{35}
\newcommand{\imgHighOptimalPten}{17}
\newcommand{\imgHighOptimalMin}{3}
\newcommand{\imgHighOptimalCount}{203}
%
% generated by FR3_Characterization/make_tx_conditions.py -- do not edit
\providecommand{\attn}[1]{\textcolor{red}{\textbf{#1}}}
\newcommand{\lmxDriveRule}{LMX2820 drive level 0 for LO frequencies up to 13\,GHz, 2 from above 13 to below 20\,GHz, and 4 from 20\,GHz up}
\newcommand{\condGainMinC}{80}
\newcommand{\condGainMaxC}{85}
\newcommand{\condFlatMinC}{80}
\newcommand{\condFlatMaxC}{85}
\newcommand{\condSbMinC}{69}
\newcommand{\condSbMaxC}{85}
\newcommand{\condCompMinC}{80}
\newcommand{\condCompMaxC}{86}
\newcommand{\condLoHotMinC}{79}
\newcommand{\condLoHotMaxC}{85}
\newcommand{\condTwoLoMinC}{79}
\newcommand{\condTwoLoMaxC}{85}
% drive rule, test temperatures, \attn
% generated by FR3_Characterization/plot_datasheet_sideband.py from sideband_tx_20261009 -- do not edit
\newcommand{\sbTxMinC}{69}
\newcommand{\sbTxMaxC}{85}
\newcommand{\sbUpperBypassMedian}{15}
\newcommand{\sbUpperBypassPten}{8}
\newcommand{\sbUpperBypassMin}{2}
\newcommand{\sbUpperBypassCount}{209}
\newcommand{\sbUpperOptimalMedian}{52}
\newcommand{\sbUpperOptimalPten}{40}
\newcommand{\sbUpperOptimalMin}{20}
\newcommand{\sbUpperOptimalCount}{209}
\newcommand{\sbLowerBypassMedian}{17}
\newcommand{\sbLowerBypassPten}{11}
\newcommand{\sbLowerBypassMin}{3}
\newcommand{\sbLowerBypassCount}{203}
\newcommand{\sbLowerOptimalMedian}{49}
\newcommand{\sbLowerOptimalPten}{37}
\newcommand{\sbLowerOptimalMin}{28}
\newcommand{\sbLowerOptimalCount}{203}

% generated by FR3_Characterization/make_sec10_1_numbers.py from results\unit_005\2026-10-09\lo_cal_grid_100MHz_20261009_044140_sec10_1_2lo_hot_table.csv -- do not edit
\newcommand{\twoLoLoMin}{7}
\newcommand{\twoLoLoMax}{13.2}
\newcommand{\twoLoLoCount}{63}
\newcommand{\twoLoReadCount}{630}
\newcommand{\twoLoMedian}{−38.7}
\newcommand{\twoLoMax}{−12.3}
\newcommand{\twoLoMaxLo}{8.1}
\newcommand{\twoLoMaxFilter}{0x15}
\newcommand{\twoLoAboveTwenty}{5}
\newcommand{\twoLoAboveLoMin}{7.8}
\newcommand{\twoLoAboveLoMax}{8.3}
\newcommand{\twoLoMinC}{79}
\newcommand{\twoLoMaxC}{85}
% the 2xLO row (all filter settings)
% generated by FR3_Characterization/make_sec6_cal_numbers.py -- do not edit
% hot: results\unit_005\2026-10-09\lo_cal_grid_100MHz_20261009_022928_full_cal_hot_drive_rule_table.csv
% cold: results\unit_005\2026-10-09\lo_cal_grid_100MHz_20261009_112919_full_cal_cold_drive_rule_table.csv
\providecommand{\attn}[1]{\textcolor{red}{\textbf{#1}}}
\providecommand{\ph}[1]{\textcolor{red}{\textbf{[#1]}}}
\newcommand{\calLoCount}{156}
\newcommand{\calReadCount}{1591}
\newcommand{\calThreshold}{−20}
\newcommand{\calPassPct}{85}
\newcommand{\calColdMinC}{57}
\newcommand{\calColdMaxC}{63}
\newcommand{\calHotMinC}{79}
\newcommand{\calHotMaxC}{85}
\newcommand{\calColdPct}{88.9}
\newcommand{\calHotPct}{90.4}
\newcommand{\calColdMedian}{−29}
\newcommand{\calHotMedian}{−30}
\newcommand{\calColdWorst}{−10}
\newcommand{\calHotWorst}{−11}
\newcommand{\calRefMedian}{−65}
\newcommand{\calMinutesPerTemp}{57}
% the LO section's numbers (also quoted in the notes below)
\raggedbottom% tables follow each other without stretched gaps
\newlength{\summarywidth}\setlength{\summarywidth}{0.84\linewidth}% all five tables share this width

The headline numbers of this datasheet in one place, for the Pi-Radio FR3-1ch unit: a circuit board (PCB revision
1.4) enclosed in a CNC-milled aluminum case, measured on Unit~\#5. Every
measured value below is taken from the section given in the last column, where it is explained in full.

\par\medskip\noindent\begin{minipage}{\linewidth}% caption and table stay on one page
\begin{center}
\captionof{table}{General.}\label{tab:summary-general}
\vspace{0.4em}
\footnotesize
\begin{tabularx}{\summarywidth}{|l|X|}
\hline
FR3 (RF) frequency range & 6--24\,GHz \\
FR1 (IF) frequency range & 0.7--4.0\,GHz \\
Instantaneous bandwidth & Up to 500\,MHz \\
LO (internal LMX2820) & 7--22.6\,GHz\textsuperscript{*}; or an external LO via the ``LO In'' connector \\
Reference & Internal 10\,MHz TCXO; or an external reference via the ``REFCLK In'' connector \\
Power & 7.5\,W from 5\,V over USB-C (supply must source 1.5\,A) \\
Control & USB-C virtual serial port, plain ASCII commands; Pi-Radio GUI \\
Intended environment & Designed for in-lab use \\
\hline
\end{tabularx}
\par\vspace{0.2em}{\scriptsize\textsuperscript{*}The LO can work as low as 6\,GHz, with minor performance degradation.\par}
\end{center}
\end{minipage}\par

\par\medskip\noindent\begin{minipage}{\linewidth}% caption and table stay on one page
\begin{center}
\captionof{table}{Absolute maximum ratings and port requirements.}\label{tab:summary-absmax}
\vspace{0.4em}
\footnotesize
\begin{tabularx}{\summarywidth}{|l|X|}
\hline
``RX-side FR3 RF In'' connector & $-$5\,dBm maximum input \\
``TX-side FR1 IF In'' connector & 0\,dBm maximum input \\
``REFCLK In'' connector & 10\,MHz, 3.3\,V CMOS clock \\
``LO In'' connector & CW sine wave, +5\,dBm, 6--22.6\,GHz \\
LMX2820 die temperature & Below 100\,\textdegree C for correct operation; beyond this limit, operation is unsafe and
not characterized. Read it in the GUI, or with the \texttt{LMX TEMP?} command. \\
ESD & \textbf{Static-sensitive device.} Some parts inside the unit can be damaged by a static discharge of under 150\,V,
far less than a person can carry without feeling it. Ground yourself, and touch each cable's outer shell to ground,
before connecting anything to any connector. \\
\hline
\end{tabularx}
\end{center}
\end{minipage}\par

\par\medskip\noindent\begin{minipage}{\linewidth}% caption and table stay on one page
\begin{center}
\captionof{table}{Transmit (FR1 $\rightarrow$ FR3). Measured across all (FR1, FR3) combinations unless noted.}\label{tab:summary-tx}
\vspace{0.4em}
\footnotesize
% Generated by FR3_Characterization/make_summary_table.py -- do not edit by hand.
\begin{tabularx}{\summarywidth}{|l|X|ccc|l|c|}
\hline
\textbf{Parameter} & \textbf{Condition} & \textbf{Min} & \textbf{Typ} & \textbf{Max} & \textbf{Unit} & \textbf{Sec.} \\
\hline
\multicolumn{7}{|l|}{\textit{Conversion gain and flatness}} \\
\hline
Conversion gain & upper sideband, filter optimal & 18.1 & 28.1 & 43.7 & dB & \ref{sec:tx-gain-upper} \\
Conversion gain & upper sideband, filter bypassed & 27.2 & 36.0 & 48.6 & dB & \ref{sec:tx-gain-upper} \\
Conversion gain & lower sideband, filter optimal & 17.0 & 30.8 & 41.9 & dB & \ref{sec:tx-gain-lower} \\
Conversion gain & lower sideband, filter bypassed & 29.5 & 37.1 & 49.3 & dB & \ref{sec:tx-gain-lower} \\
Gain flatness (RMS) & upper sideband, filter optimal & 0.17 & 1.51 & 3.29 & dB & \ref{sec:tx-flatness} \\
Gain flatness (RMS) & upper sideband, filter bypassed & 0.20 & 1.55 & 3.63 & dB & \ref{sec:tx-flatness} \\
Gain flatness (RMS) & lower sideband, filter optimal & 0.16 & 1.26 & 3.07 & dB & \ref{sec:tx-flatness} \\
Gain flatness (RMS) & lower sideband, filter bypassed & 0.19 & 1.21 & 2.74 & dB & \ref{sec:tx-flatness} \\
\hline
\multicolumn{7}{|l|}{\textit{Compression}} \\
\hline
Output P1dB & upper sideband, 9 test points & 8.1 & 17.1 & 18.2 & dBm & \ref{sec:compression-upper} \\
Input P1dB & upper sideband, 9 test points & −24.3 & −13.2 & −7.0 & dBm & \ref{sec:compression-upper} \\
Output P1dB & lower sideband, 9 test points & 4.9 & 15.8 & 18.4 & dBm & \ref{sec:compression-lower} \\
Input P1dB & lower sideband, 9 test points & −21.1 & −11.5 & −5.6 & dBm & \ref{sec:compression-lower} \\
\hline
\multicolumn{7}{|l|}{\textit{Sideband suppression}} \\
\hline
Sideband suppression & upper desired, filter optimal & 40\raisebox{0.55ex}{\scriptsize a} & 52 &  & dB & \ref{sec:sideband-suppression} \\
Sideband suppression & upper desired, filter bypassed & 8\raisebox{0.55ex}{\scriptsize a} & 15 &  & dB & \ref{sec:sideband-suppression} \\
Sideband suppression & lower desired, filter optimal & 37\raisebox{0.55ex}{\scriptsize a} & 49 &  & dB & \ref{sec:sideband-suppression} \\
Sideband suppression & lower desired, filter bypassed & 11\raisebox{0.55ex}{\scriptsize a} & 17 &  & dB & \ref{sec:sideband-suppression} \\
\hline
\multicolumn{7}{|l|}{\textit{LO leakage}} \\
\hline
LO leakage, other filter settings & calibrated, hot (LMX 79--85\,\textdegree C) &  & −30\raisebox{0.55ex}{\scriptsize b} & −11 & dBm & \ref{sec:lo-suppression} \\
LO leakage below −20\,dBm & calibrated, hot, all (LO, filter) readings &  & 90.4 &  & \% & \ref{sec:lo-suppression} \\
\hline
\multicolumn{7}{|l|}{\textit{Spurious outputs}} \\
\hline
2×LO spur & LO 7--13.2\,GHz, all filter settings, no input\raisebox{0.55ex}{\scriptsize c} &  & −38.7 & −12.3 & dBm & \ref{sec:spur-tx-2lo} \\
\hline
\end{tabularx}

\end{center}
\end{minipage}\par

\par\medskip\noindent\begin{minipage}{\linewidth}% caption and table stay on one page
\begin{center}
\captionof{table}{Receive (FR3 $\rightarrow$ FR1). Measured across all (FR1, FR3) combinations unless noted.}\label{tab:summary-rx}
\vspace{0.4em}
\footnotesize
% Generated by FR3_Characterization/make_summary_table.py -- do not edit by hand.
\begin{tabularx}{\summarywidth}{|l|X|ccc|l|c|}
\hline
\textbf{Parameter} & \textbf{Condition} & \textbf{Min} & \textbf{Typ} & \textbf{Max} & \textbf{Unit} & \textbf{Sec.} \\
\hline
\multicolumn{7}{|l|}{\textit{Conversion gain and flatness}} \\
\hline
Conversion gain & low-side LO, filter optimal & 27.5 & 37.2 & 53.0 & dB & \ref{sec:rx-gain-lowside} \\
Conversion gain & low-side LO, filter bypassed & 37.7 & 45.9 & 56.9 & dB & \ref{sec:rx-gain-lowside} \\
Conversion gain & high-side LO, filter optimal & 27.2 & 39.6 & 51.2 & dB & \ref{sec:rx-gain-highside} \\
Conversion gain & high-side LO, filter bypassed & 36.7 & 46.0 & 57.5 & dB & \ref{sec:rx-gain-highside} \\
Gain flatness (RMS) & low-side LO, filter optimal & 0.39 & 0.95 & 2.24 & dB & \ref{sec:rx-flatness} \\
Gain flatness (RMS) & low-side LO, filter bypassed & 0.46 & 1.03 & 2.63 & dB & \ref{sec:rx-flatness} \\
Gain flatness (RMS) & high-side LO, filter optimal & 0.38 & 0.90 & 3.19 & dB & \ref{sec:rx-flatness} \\
Gain flatness (RMS) & high-side LO, filter bypassed & 0.37 & 0.95 & 2.89 & dB & \ref{sec:rx-flatness} \\
\hline
\multicolumn{7}{|l|}{\textit{Compression}} \\
\hline
Input P1dB & low-side LO, 8 test points & −32.6 & −23.9 & −12.3 & dBm & \ref{sec:rx-compression-lowside} \\
Output P1dB & low-side LO, 8 test points & 13.2 & 14.3 & 18.4 & dBm & \ref{sec:rx-compression-lowside} \\
Input P1dB & high-side LO, 9 test points & −28.3 & −23.2 & −13.7 & dBm & \ref{sec:rx-compression-highside} \\
Output P1dB & high-side LO, 9 test points & 13.2 & 14.5 & 18.3 & dBm & \ref{sec:rx-compression-highside} \\
\hline
\multicolumn{7}{|l|}{\textit{Image rejection}} \\
\hline
Image rejection & low-side LO, filter optimal & 24\raisebox{0.55ex}{\scriptsize a} & 51 &  & dB & \ref{sec:rx-image-rejection} \\
Image rejection & high-side LO, filter optimal & 17\raisebox{0.55ex}{\scriptsize a} & 35 &  & dB & \ref{sec:rx-image-rejection} \\
\hline
\multicolumn{7}{|l|}{\textit{Noise figure and sensitivity}} \\
\hline
Noise figure (SSB) & low-side LO, filter optimal & 5.2 & 12.2 & 22.9 & dB & \ref{sec:rx-noise-figure} \\
Noise figure (SSB) & high-side LO, filter optimal & 5.6 & 11.0 & 20.5 & dB & \ref{sec:rx-noise-figure} \\
Noise floor at RF input & median NF, filter optimal, 100\,MHz bandwidth &  & −82.5 &  & dBm & \ref{sec:rx-noise-figure} \\
\hline
\multicolumn{7}{|l|}{\textit{Spurious outputs}} \\
\hline
LO leakage at RF input & LO 7--22.5\,GHz, TX LO nulled\raisebox{0.55ex}{\scriptsize d} &  & −59.6 & −49.3 & dBm & \ref{sec:spur-rx-lo} \\
2×LO leakage at RF input & LO 7--13.2\,GHz, TX LO nulled\raisebox{0.55ex}{\scriptsize d} &  & −71.1 & −58.8 & dBm & \ref{sec:spur-rx-lo} \\
\hline
\end{tabularx}

\end{center}
\end{minipage}\par

\par\medskip\noindent\begin{minipage}{\linewidth}% caption and table stay on one page
\begin{center}
\captionof{table}{LO and phase noise.}\label{tab:summary-lo}
\vspace{0.4em}
\footnotesize
% Generated by FR3_Characterization/make_summary_table.py -- do not edit by hand.
\begin{tabularx}{\summarywidth}{|l|X|ccc|l|c|}
\hline
\textbf{Parameter} & \textbf{Condition} & \textbf{Min} & \textbf{Typ} & \textbf{Max} & \textbf{Unit} & \textbf{Sec.} \\
\hline
RMS phase error & 100\,Hz--1\,MHz, internal LO, 8--22\,GHz & 0.82 &  & 2.13 & degrees & \ref{sec:phase-noise} \\
Phase noise & 100\,kHz offset, internal LO at 10\,GHz &  & −95.1 &  & dBc/Hz & \ref{sec:phase-noise} \\
RMS phase error & Pi-Radio OctoLO (external LO) at 10\,GHz &  & 0.24 &  & degrees & \ref{sec:phase-noise} \\
\hline
\end{tabularx}

\end{center}
\end{minipage}\par

\noindent{\footnotesize\textbf{Notes:} \textbf{Typ} is the median over every measured point within the design spec
(FR3~$\le$~24\,GHz; for conversion gain and gain flatness, whose 500\,MHz window at FR1~=~4\,GHz runs past the IF spec, also FR1
centre~$\le$~3.5\,GHz), and \textbf{Min}/\textbf{Max} are the extremes over the same points;
points outside the design spec are shown in the sections but left out here. All TX values were measured with the unit hot (no fan, LMX2820 die at about \condGainMinC--\condCompMaxC\textdegree C), as the worst case, with \lmxDriveRule. Conversion gain is the average over a
500\,MHz window at each point, measured at a $-$40\,dBm input. Gain flatness is the RMS deviation of the gain across that same 500\,MHz window. P1dB is measured at the representative test points of
Sections~\ref{sec:compression-upper}, \ref{sec:compression-lower}, \ref{sec:rx-compression-lowside}, and \ref{sec:rx-compression-highside} only. Noise figure is given with the AM3153 filter bypassed and with the ``optimal'' filter engaged
(Section~\ref{sec:rx-noise-figure}).
\par\noindent\raisebox{0.55ex}{\scriptsize a}\,10th percentile: 90\% of measured points are at or above this value (TX: computed from the measured points; RX: read from the distribution figure, to within about 1\,dB).
\par\noindent\raisebox{0.55ex}{\scriptsize b}\,Median, and worst, LO leakage at the ``TX-side FR3 RF Out'' connector with the stored calibration applied, over every filter setting usable at each of the \calLoCount{} calibrated LO frequencies (7.0--22.5\,GHz in 100\,MHz steps), excluding the setting the null was found in; the full sweep and its distribution are in Section~\ref{sec:lo-suppression}.
\par\noindent\raisebox{0.55ex}{\scriptsize c}\,Median, and worst, 2×LO power at the ``TX-side FR3 RF Out'' connector over all \twoLoReadCount{} readings: at each of \twoLoLoCount{} LO frequencies (\twoLoLoMin--\twoLoLoMax\,GHz in 100\,MHz steps), the LO leakage was nulled with the reference filter and the 2×LO was then read in every filter setting usable at that LO, with no input signal. The worst reading is at LO~=~\twoLoMaxLo\,GHz with filter \twoLoMaxFilter; \twoLoAboveTwenty{} of the \twoLoReadCount{} readings are above $-$20\,dBm. Above LO~=~13.25\,GHz the spur lies above 26.5\,GHz and was not measured (Section~\ref{sec:spur-tx-2lo}).
\par\noindent\raisebox{0.55ex}{\scriptsize d}\,Power leaking out of the ``RX-side FR3 RF In'' connector (radiated if an antenna is connected), with the TX side in normal use; 2×LO above \spurTwoLoLimit\,GHz not measured (Section~\ref{sec:spur-rx-lo}). \par}% \par inside the group: line spacing follows \footnotesize
\flushbottom
